\documentclass[11pt]{article}
\usepackage[margin=1in]{geometry}
\usepackage[T1]{fontenc}
\usepackage[utf8]{inputenc}
\title{Lynjax v1.0-rc1 Manual\\Runnable Test Candidate Package}
\author{Hermes Agent para Alejandro}
\date{2026-06-15}
\begin{document}
\maketitle

\section*{Estado}
Lynjax v1.0-rc1 es un candidato tecnico de prueba, no la publicacion final v1.0. El objetivo es validar ejecucion local y base virtual/container-lab sin credenciales reales, sin escaneo externo y sin cambios de host.

\section{Alcance}
\begin{itemize}
  \item Frontend React/Vite bilingue ES/EN con shell de plataforma, sidebar y topbar.
  \item Backend FastAPI con health, metadata y endpoint /api/v1/assessments/connectivity-demo.
  \item Flujo demo assessment a evidencia a reporte Markdown.
  \item Docker Compose beta para WSL2/Ubuntu/VM/CI.
  \item Topologia Containerlab sanitizada y validada estaticamente.
\end{itemize}

\section{Capturas generadas}
Las capturas del frontend y backend fueron generadas y guardadas en:
\begin{itemize}
  \item docs/manual/assets/screenshots/lynjax-v1.0-rc1-frontend-overview.png
  \item docs/manual/assets/screenshots/lynjax-v1.0-rc1-backend-openapi.png
\end{itemize}

\section{Comandos locales}
\begin{verbatim}
cd /c/Users/nesal/Documents/001_Programas/netvault-rebrand-lab/lynjax
python -m pytest backend/tests -v
npm --prefix frontend install
npm --prefix frontend run build
bash -n scripts/*.sh virtualization/run-beta-compose.sh
bash scripts/lab_validate.sh
bash scripts/dev-start.sh
bash scripts/smoke-local.sh
bash scripts/dev-stop.sh
\end{verbatim}

\section{Comandos WSL2/Ubuntu con Docker}
\begin{verbatim}
cd /mnt/c/Users/nesal/Documents/001_Programas/netvault-rebrand-lab/lynjax
. /home/nstalej/.nvm/nvm.sh
nvm use 24
backend/.venv-wsl/bin/python -m pytest backend/tests -v
npm --prefix frontend install
npm --prefix frontend run build
bash scripts/lab_validate.sh
cd virtualization
bash run-beta-compose.sh config
bash run-beta-compose.sh up-detached
curl -fsS http://127.0.0.1:8000/health
curl -fsS http://127.0.0.1:5173/ >/dev/null
curl -fsS http://127.0.0.1:18080/ >/dev/null
curl -fsS http://127.0.0.1:18081/metadata.json >/dev/null
bash run-beta-compose.sh down -v
\end{verbatim}

\section{Matriz de validacion}
\begin{itemize}
  \item Backend pytest: PASS, 5 tests en Windows y WSL2.
  \item Frontend build: PASS, Vite build OK tras npm install.
  \item Bash syntax: PASS.
  \item Lab validate: PASS; Compose config OK en WSL2.
  \item Local smoke: PASS.
  \item Docker Compose smoke: PASS para backend, frontend, target web y metadata.
  \item Containerlab runtime: SKIP; containerlab no instalado.
  \item NPM audit: WARN; 2 high por Vite/esbuild, upgrade mayor diferido.
\end{itemize}

\section{Riesgos y limites}
\begin{itemize}
  \item No publicar v1.0 final sin aprobacion explicita.
  \item Resolver o aceptar formalmente los hallazgos NPM antes de produccion.
  \item Instalar/probar Containerlab solo dentro de WSL2/Ubuntu/VM aprobada.
  \item Mantener alcance sandbox: sin credenciales reales ni redes externas.
\end{itemize}

\section{Proximos pasos manuales}
\begin{enumerate}
  \item Revisar navegacion y selector ES/EN.
  \item Ejecutar la demo segura desde frontend y Swagger.
  \item Repetir smoke WSL2/Compose.
  \item Decidir si se etiqueta final v1.0 o si se abre rama de hardening Vite/Containerlab.
\end{enumerate}

\end{document}
